\documentclass[11pt,respuestas,a4]{aleph-examen}
% \documentclass[10pt,a5]{aleph-examen}

% -- Paquetes extra
\usepackage{aleph-comandos}
\usepackage{booktabs}
\usepackage{multicol}
\usepackage{systeme}

% -- Datos
\institucion{Facultad de Ciencias Exactas, Naturales y Ambientales}
\carrera{Catalogo STEM}
\asignatura{Algebra lineal}
\tema{Examen no. 2: Espacios vectoriales}
\autor{Andres Merino}
\fecha{Periodo 2026-1}

\logouno[0.16\textwidth]{Logos/logoPUCE_04_ac}
\definecolor{colortext}{HTML}{1957d1}
\definecolor{colordef}{HTML}{1957d1}
\fuente{montserrat}

\begin{document}

\encabezado

\vspace{-5mm}\small
\section*{Indicaciones}
\begin{itemize}[leftmargin=*]
\item
    En esta actividad se evalua si el estudiante (Citerios 3.1) resuelve operaciones en espacios vectoriales y aplicaciones lineales, incluyendo el cálculo de bases, núcleo, imagen, así como valores y vectores propios.
\item
    Los ejercicios deben ser resueltos a mano, sin ayuda de resumenes o asistentes computacionales.
\item
    Las resoluciones deben entregarse en hojas de forma fisica al docente.
\item
    Se encuentra prohibido el uso de cualquier otra fuente de informacion durante todo el examen.
\item
    En caso de considerar que existe un error en la pregunta o que esta se encuentra mal planteada, se debe indicar cual es el error y justificarlo.
\item
    Todas las soluciones deben estar correctamente redactadas y explicadas. \puntaje{1}
\end{itemize}

\section*{Ejercicios}

\begin{preguntas}

%%%%%%%%%%%%%%%%%%%%%%%%%%%%%%%%%%%%%%%%
%% Base
%%%%%%%%%%%%%%%%%%%%%%%%%%%%%%%%%%%%%%%%
\item
    Responda las siguientes preguntas, justificando sus respuestas.
    \begin{enumerate}[leftmargin=*]
    \item 
        Considere $B$ un subconjunto de $\R^{5}$. Si $B$ tiene $4$ elementos, ¿puede ser $B$ un conjunto linealmente independiente?, ¿puede ser $B$ un conjunto generador de $E$?\puntaje{1}
    \item 
        Considere $B$ un subconjunto de $\R^{2\times 2}$. Si $B$ tiene $5$ elementos, ¿puede ser $B$ un conjunto linealmente independiente?, ¿puede ser $B$ un conjunto generador de $E$?\puntaje{1}
    \item 
        En $\R^{2\times 3}$, ¿puede un conjunto de $5$ elementos ser base?, ¿puede un conjunto de $7$ elementos ser base?\puntaje{1}
    \item 
        En cualquier espacio, ¿por qué un conjunto que contenga el vector nulo no puede ser base del espacio?\puntaje{1}
    \end{enumerate}

\begin{respuesta}[Respuesta]\hspace{0pt}
    \begin{enumerate}
    \item 
        Como $|B|=4$ y $\dim(\R^{5}) = 5$, tenemos que $|B| < \dim(\R^{5})$, por lo tanto, el conjunto no puede ser un conjunto generador. Por otro lado, sí puede ser un conjunto linealmente independiente. 
    \item 
        Como $|B|=5$ y $\dim(\R^{2\times 2}) = 4$, tenemos que $|B| > \dim(\R^{2\times 2})$, por lo tanto, el conjunto no puede ser un conjunto linealmente independiente. Por otro lado, sí puede ser un conjunto generador.
    \item
        Como $\dim(\R^{2\times 3}) = 6$, para que un conjunto sea base, debe tener $6$ elementos, por lo tanto, la respuesta a ambas preguntas es no.
    \item 
        Si el conjunto contiene el vector nulo, entonces el conjunto no es linealmente independiente, por lo tanto, no puede ser base. \qedhere
    \end{enumerate}
\end{respuesta}

%%%%%%%%%%%%%%%%%%%%%%%%%%%%%%%%%%%%%%%%
%% Conjunto generado
%%%%%%%%%%%%%%%%%%%%%%%%%%%%%%%%%%%%%%%%
\item
    En $\R^4$, se considera el conjunto
    \[
        S=\big\{
            (1,0,1,1),
            (0,1,1,-1)
        \big\}.
    \]
    \begin{enumerate}
    \item
        Determine el conjunto generado por $S$ y escribalo en funcion de sus restricciones.\puntaje{8}
    \item
        Sin resolver un nuevo sistema completo, determine si $(3,1,4,2)$ y $(1,2,3,0)$ pertenecen a $\gen(S)$.\puntaje{2}
    \end{enumerate}

\begin{respuesta}\hspace{0pt}
\begin{enumerate}
\item
    Si $(w,x,y,z)\in\gen(S)$, entonces existen $a,b\in\R$ tales que
    \[
        a(1,0,1,1)+b(0,1,1,-1)=(w,x,y,z).
    \]
    Esto equivale al sistema
    \[
        \systeme[ab]{
            a=w,
            b=x,
            a+b=y,
            a-b=z
        }.
    \]
    Su matriz ampliada es
    \[
        \begin{pmatrix}
            1 & 0 & | & w\\
            0 & 1 & | & x\\
            1 & 1 & | & y\\
            1 & -1 & | & z
        \end{pmatrix}
        \sim
        \begin{pmatrix}
            1 & 0 & | & w\\
            0 & 1 & | & x\\
            0 & 0 & | & y-w-x\\
            0 & 0 & | & z-w+x
        \end{pmatrix}.
    \]
    Por lo tanto, el sistema tiene solucion si y solo si
    \[
        y=w+x
        \qquad\text{y}\qquad
        z=w-x.
    \]
    Por lo tanto,
    \[
        \gen(S)=
        \big\{(w,x,y,z)\in\R^4:\ y=w+x\ \land\ z=w-x\big\}.
    \]
\item
    Para $(3,1,4,2)$, se cumple
    \[
        4=3+1
        \qquad\text{y}\qquad
        2=3-1.
    \]
    Por lo tanto, $(3,1,4,2)\in\gen(S)$.

    Para $(1,2,3,0)$, se cumple $3=1+2$, pero
    \[
        0\neq 1-2.
    \]
    Por lo tanto, $(1,2,3,0)\notin\gen(S)$.\qedhere
\end{enumerate}
\end{respuesta}

%%%%%%%%%%%%%%%%%%%%%%%%%%%%%%%%%%%%%%%%
%% Independencia lineal
%%%%%%%%%%%%%%%%%%%%%%%%%%%%%%%%%%%%%%%%
\item
    Sea $\alpha\in\R$. En $\R^4$, se considera el conjunto
    \[
        B_\alpha=\big\{
            (1,0,0,1),
            (0,1,0,1),
            (2,2,\alpha,4)
        \big\}.
    \]
    \begin{enumerate}
    \item
        Determine para que valores de $\alpha$ el conjunto $B_\alpha$ es linealmente independiente.\puntaje{8}
    \end{enumerate}

\begin{respuesta}\hspace{0pt}
\begin{enumerate}
\item
    Para analizar si $B_\alpha$ es linealmente independiente, planteamos
    \[
        c_1(1,0,0,1)+c_2(0,1,0,1)+c_3(2,2,\alpha,4)=(0,0,0,0).
    \]
    Esto equivale al sistema homogeneo
    \[
        \systeme[c_1c_2c_3]{
            c_1+2c_3=0,
            c_2+2c_3=0,
            \alpha c_3=0,
            c_1+c_2+4c_3=0
        }.
    \]
    La matriz ampliada asociada es
    \[
        \begin{pmatrix}
            1 & 0 & 2 & | & 0\\
            0 & 1 & 2 & | & 0\\
            0 & 0 & \alpha & | & 0\\
            1 & 1 & 4 & | & 0
        \end{pmatrix}
        \sim
        \begin{pmatrix}
            1 & 0 & 2 & | & 0\\
            0 & 1 & 2 & | & 0\\
            0 & 0 & \alpha & | & 0\\
            0 & 0 & 0 & | & 0
        \end{pmatrix}.
    \]
    Si $\alpha\neq 0$, entonces el sistema homogeneo tiene solamente la solucion trivial y $B_\alpha$ es linealmente independiente.

    Si $\alpha=0$, el sistema tiene infinitas soluciones y por ende el conjunto es linealmente dependiente.

    En consecuencia,
    \[
        B_\alpha \text{ es linealmente independiente si y solo si } \alpha\neq 0.
        \qedhere
    \]
\end{enumerate}
\end{respuesta}

%%%%%%%%%%%%%%%%%%%%%%%%%%%%%%%%%%%%%%%%
%% Aplicaciones lineales
%%%%%%%%%%%%%%%%%%%%%%%%%%%%%%%%%%%%%%%%
\item
    Sea $T:\R^4\to\R^3$ la aplicacion lineal definida por
    \[
        T(x,y,z,w)=(x+y,\ y+z,\ x+2y+z).
    \]
    \begin{enumerate}
    \item
        Sin calcular el nucleo ni la imagen, indique si $T$ podria ser inyectiva y si podria ser sobreyectiva.\puntaje{2}
    \item
        Determine el nucleo de $T$ y una base para este.\puntaje{8}
    \item
        Determine la imagen de $T$ y una base para esta.\puntaje{8}
    \item
        Con lo anterior, determine si $T$ es inyectiva, sobreyectiva o isomorfismo.\puntaje{2}
    \end{enumerate}

\begin{respuesta}\hspace{0pt}
\begin{enumerate}
\item
    Como $\dim(\R^4)=4$ y $\dim(\R^3)=3$, una aplicacion lineal de $\R^4$ en $\R^3$ no puede ser inyectiva. En cambio, si podria ser sobreyectiva, pues el espacio de llegada tiene menor dimension que el espacio de partida.
\item
    Para calcular el nucleo, resolvemos
    \[
        T(x,y,z,w)=(0,0,0),
    \]
    es decir,
    \[
        (x+y,\ y+z,\ x+2y+z)=(0,0,0).
    \]
    Entonces,
    \[
        \systeme[xyzw]{
            x+y=0{,},
            y+z=0{,},
            x+2y+z=0{.}
        }
    \]
    La matriz ampliada del sistema es
    \[
        \begin{pmatrix}
            1 & 1 & 0 & 0 & | & 0\\
            0 & 1 & 1 & 0 & | & 0\\
            1 & 2 & 1 & 0 & | & 0
        \end{pmatrix}
        \sim
        \begin{pmatrix}
            1 & 0 & -1 & 0 & | & 0\\
            0 & 1 & 1 & 0 & | & 0\\
            0 & 0 & 0 & 0 & | & 0
        \end{pmatrix}.
    \]
    Por lo tanto, el sistema equivalente es
    \[
        x=-y
        \qquad\text{y}\qquad
        z=-y.
    \]
    Tomando $y=t$ y $w=s$, se tiene
    \[
        (x,y,z,w)=(-t,t,-t,s)
        =
        t(-1,1,-1,0)+s(0,0,0,1).
    \]
    Por lo tanto,
    \[
        \ker(T)=
        \gen\big\{(-1,1,-1,0),(0,0,0,1)\big\}.
    \]
    Una base para $\ker(T)$ es
    \[
        \big\{(-1,1,-1,0),(0,0,0,1)\big\}.
    \]
\item
    Tomamos $(a,b,c)\in\R^3$ y buscamos
    $r,s,t\in\R$ tales que
    \[
        r(1,0,1)+s(1,1,2)+t(0,1,1)=(a,b,c).
    \]
    Esto equivale al sistema
    \[
        \systeme[rst]{
            r+s=a,
            s+t=b,
            r+2s+t=c
        }.
    \]
    Su matriz ampliada es
    \[
        \begin{pmatrix}
            1 & 1 & 0 & | & a\\
            0 & 1 & 1 & | & b\\
            1 & 2 & 1 & | & c
        \end{pmatrix}
        \sim
        \begin{pmatrix}
            1 & 0 & -1 & | & a-b\\
            0 & 1 & 1 & | & b\\
            0 & 0 & 0 & | & c-a-b
        \end{pmatrix}.
    \]
    Entonces, el sistema tiene solucion si y solo si $c=a+b$. Por lo tanto,
    \[
        \img(T)=\big\{(a,b,c)\in\R^3:\ c=a+b\big\}.
    \]
    Además, como
    \[
        (a,b,c) = (a,b,a+b) = a(1,0,1)+b(0,1,1),
    \]
    se tiene que
    \[
        \img(T)=
        \gen\big\{(1,0,1),(0,1,1)\big\}.
    \]
    Así, una base para la imagen es
    \[
        \big\{(1,0,1),(0,1,1)\big\}.
    \]
\item
    Como $\ker(T)\neq\{(0,0,0,0)\}$, la aplicacion no es inyectiva. Como
    \[
        \img(T)=\big\{(a,b,c)\in\R^3:\ c=a+b\big\}\neq\R^3,
    \]
    la aplicacion no es sobreyectiva. En consecuencia, $T$ no es un isomorfismo.\qedhere
\end{enumerate}
\end{respuesta}

%%%%%%%%%%%%%%%%%%%%%%%%%%%%%%%%%%%%%%%%
%% Aplicaciones lineales
%%%%%%%%%%%%%%%%%%%%%%%%%%%%%%%%%%%%%%%%
\item
    Considere la siguiente aplicación lineal:
    \[
        \funcion{T}{\R^4}{\R^{2\times 2}}
            {(a,b,c,d)}{\displaystyle\begin{pmatrix} a+b && b+c\\ c+d && d+a \end{pmatrix};}
    \]
    además, considere las siguientes bases de $\R^4$ y $\R^{2\times 2}$, respectivamente:
    \[
        B_1=\{ (1,2,0,0),\ (0,1,2,0),\ (0,0,1,2),\ (0,0,0,2)\}
    \]
    y
    \[
        B_2=\left\{
            \begin{pmatrix} 1 & 0 \\ 0 & 0 \end{pmatrix},\ 
            \begin{pmatrix} 0 & 1 \\ 0 & 0 \end{pmatrix},\ 
            \begin{pmatrix} 0 & 0 \\ 1 & 0 \end{pmatrix},\ 
            \begin{pmatrix} 0 & 0 \\ 0 & 1 \end{pmatrix}
        \right\}.
    \]
    \begin{enumerate}
    \item
        Evalúe la aplicación $T$ en cada elemento de $B_1$. \puntaje{2}
    \item
        Determine el vector de coordenadas en la base $B_2$ de cada uno de los resultados del punto anterior. \puntaje{2}
    \item
        Determine la matriz asociada a la aplicación lineal $T$, es decir, determine $[T]$. \puntaje{1}
    \item
        ¿Para qué sirve $[T]$? \puntaje{2}
    \end{enumerate}

\begin{respuesta}\hspace{0pt}
    \begin{enumerate}
    \item 
        Tenemos que
        \begin{align*}
            T(1,2,0,0) & = \begin{pmatrix}
                1+2 & 2+0 \\
                0+0 & 0+1
            \end{pmatrix} = \begin{pmatrix}
                3 & 2 \\
                0 & 1
            \end{pmatrix},\\[2mm]
            T(0,1,2,0) & = \begin{pmatrix}
                0+1 & 1+2 \\
                2+0 & 0+0
            \end{pmatrix} = \begin{pmatrix}
                1 & 3 \\
                2 & 0
            \end{pmatrix},\\[2mm]
            T(0,0,1,2) & = \begin{pmatrix}
                0+0 & 0+1 \\
                1+2 & 2+0
            \end{pmatrix} = \begin{pmatrix}
                0 & 1 \\
                3 & 2
            \end{pmatrix},\\[2mm]
            T(0,0,0,2) & = \begin{pmatrix}
                0+0 & 0+0 \\
                0+2 & 2+0
            \end{pmatrix} = \begin{pmatrix}
                0 & 0 \\
                2 & 2
            \end{pmatrix}.
        \end{align*}
    \item 
         Como $B_2$ es la matriz canónica, se tiene que:
        \begin{align*}
            [T(1,2,0,0)]_{B_2} &=
            \left[\begin{pmatrix}
                3 & 2 \\
                0 & 1
            \end{pmatrix}\right]_{B_2}
            =
            \begin{pmatrix}
                3 \\
                2 \\
                0 \\
                1
            \end{pmatrix}.\\[2mm]
            [T(0,1,2,0)]_{B_2} &=
            \left[\begin{pmatrix}
                1 & 3 \\
                2 & 0
            \end{pmatrix}\right]_{B_2}
            =
            \begin{pmatrix}
                1 \\
                3 \\
                2 \\
                0
            \end{pmatrix},\\[2mm]
            [T(0,0,1,2)]_{B_2} &=
            \left[\begin{pmatrix}
                0 & 1 \\
                3 & 2
            \end{pmatrix}\right]_{B_2}
            =
            \begin{pmatrix}
                0 \\
                1 \\
                3 \\
                2
            \end{pmatrix},\\[2mm]
            [T(0,0,0,2)]_{B_2} &=
            \left[\begin{pmatrix}
                0 & 0 \\
                2 & 2
            \end{pmatrix}\right]_{B_2}
            =
            \begin{pmatrix}
                0 \\
                0 \\
                2 \\
                2
            \end{pmatrix}.
        \end{align*}
    \item
        Para determinar la matriz asociada a la aplicación lineal $T$ debemos calcular la
        imagen de cada vector de la base $B_1$ y luego expresar cada imagen en la base
        $B_2$. Esos vectores de coordenadas forman las columnas de $[T]$.
        Así, se tiene que
        \[
            [T] = \begin{pmatrix}
                3 & 1 & 0 & 0 \\
                2 & 3 & 1 & 0 \\
                0 & 2 & 3 & 2 \\
                1 & 0 & 2 & 2
            \end{pmatrix}.
        \]
    \item 
        La matriz $[T]$ sirve para calcular las coordenadas de $T(v)$ en la base $B_2$ a
        partir de las coordenadas de $v$ en la base $B_1$. Es decir, para todo
        $v\in\R^4$ se cumple
        \[
            [T(v)]_{B_2}=[T]\,[v]_{B_1}.
        \]
        Por lo tanto, $[T]$ permite representar la aplicación lineal mediante una
        multiplicación matricial. \qedhere
    \end{enumerate}
\end{respuesta}

\end{preguntas}

\end{document}
